\scanpage{86}
\[
\tag{\unsure{41|42}}f_{l_i}(u)\qquad i\in I,\ u\in\Gamma
\]
où pour $\forall i\in I$, $l_i$ est un chemin dans $\Pi_1(X,x_i)$ de $x_0$ vers $x_i$, i.e. un élément de la partie ``discrète'' du torseur (41).

D'autre part, l'action de $\Gamma$ sur $\widehat\pi_1(X,x_0)$ est connue quand on connaît son action sur des générateurs $\lambda_j$ \inserted{($j\in J$)} de $\pi_1(X,x_0)$, i.e. si on connaît les
\[
\tag{43}f_{\lambda_j}(u)\quad\text{ou}\quad g_{\lambda_j}(u)\qquad j\in J,\ u\in\Gamma.
\]
\deleted{Mais cette} Les données (42)(43) suffisent donc à déterminer l'action de $\Gamma$ sur $\Pi_1(X,I)$ --- mais cette description, \deleted{à la} \deleted{\unsure{dis|dire}} s'il a l'avantage d'une économie, au moins en apparence, a le désavantage de se référer à des choix arbitraires, le pire étant le choix d'une \deleted{origine} \inserted{point-base} $x_0$ \deleted{à laquelle} \inserted{auquel} on est censé \unsure{tout} ramener --- ce qui semble artificiel, et à la longue je crois impraticable, \deleted{est} et tout particulièrement dans les situations impliquant un groupe de symétrie $\unsure{\mathcal G|G}$ comme ci-dessus, quand on a pris $I$ stable sous $\unsure{\mathcal G|G}$ ---
\margin{qu'il faut \unsure{exprimer} \unsure{linéairement|littéralement}, \unsure{resp.} \unsure{principes}, \unsure{décrit p.8!|c'est le pb!}}
\margin{85}
